\chapter{Lab 1: Membangun Aplikasi RAG Full-Stack}
\label{chap:lab1}

Dalam lab ini, kita tidak hanya menulis skrip Python, tetapi membangun aplikasi web lengkap yang bisa diakses via browser. Kita akan menggunakan **Streamlit**, framework UI paling populer untuk Data Science.

\section{Arsitektur Sistem}

Aplikasi kita akan memiliki komponen berikut:
\begin{enumerate}
    \item **Frontend:** Streamlit (Input PDF, Chat Interface).
    \item **Processing:** PyPDF2 (Extract text), LangChain (Text Splitter).
    \item **Database:** ChromaDB (Vector Store lokal).
    \item **Brain:** OpenAI GPT-3.5/4 (Generative response).
\end{enumerate}

\section{Persiapan Environment}

Buat file \texttt{requirements.txt} dan isi dengan:

\begin{codeblock}
streamlit
openai
langchain
langchain-community
chromadb
pypdf2
python-dotenv
tiktoken
\end{codeblock}

Install dependensi:
\begin{codeblock}
pip install -r requirements.txt
\end{codeblock}

\section{Kode Utama: app.py}

Berikut adalah kode lengkap untuk aplikasi \texttt{app.py}. Salin dan pelajari komentar di dalamnya.

\begin{codeblock}[language=Python]
import streamlit as st
from dotenv import load_dotenv
from PyPDF2 import PdfReader
from langchain.text_splitter import CharacterTextSplitter
from langchain_community.embeddings import OpenAIEmbeddings
from langchain_community.vectorstores import Chroma
from langchain.memory import ConversationBufferMemory
from langchain.chains import ConversationalRetrievalChain
from langchain_community.chat_models import ChatOpenAI

def get_pdf_text(pdf_docs):
    text = ""
    for pdf in pdf_docs:
        pdf_reader = PdfReader(pdf)
        for page in pdf_reader.pages:
            text += page.extract_text()
    return text

def get_text_chunks(text):
    text_splitter = CharacterTextSplitter(
        separator="\n",
        chunk_size=1000,
        chunk_overlap=200,
        length_function=len
    )
    chunks = text_splitter.split_text(text)
    return chunks

def get_vectorstore(text_chunks):
    embeddings = OpenAIEmbeddings()
    # Persist directory opsional untuk menyimpan DB di disk
    vectorstore = Chroma.from_texts(texts=text_chunks, embedding=embeddings)
    return vectorstore

def get_conversation_chain(vectorstore):
    llm = ChatOpenAI()
    memory = ConversationBufferMemory(
        memory_key='chat_history',
        return_messages=True
    )
    conversation_chain = ConversationalRetrievalChain.from_llm(
        llm=llm,
        retriever=vectorstore.as_retriever(),
        memory=memory
    )
    return conversation_chain

def handle_userinput(user_question):
    response = st.session_state.conversation({'question': user_question})
    st.session_state.chat_history = response['chat_history']

    for i, message in enumerate(st.session_state.chat_history):
        if i % 2 == 0:
            st.write(f"[User] User: {message.content}")
        else:
            st.write(f"[Bot] Bot: {message.content}")

def main():
    load_dotenv()
    st.set_page_config(page_title="Chat dengan PDF", page_icon="[Book]")
    st.header("Chat dengan PDF Anda [Book]")

    if "conversation" not in st.session_state:
        st.session_state.conversation = None
    if "chat_history" not in st.session_state:
        st.session_state.chat_history = None

    user_question = st.text_input("Ajukan pertanyaan tentang dokumen:")
    if user_question:
        handle_userinput(user_question)

    with st.sidebar:
        st.subheader("Dokumen Anda")
        pdf_docs = st.file_uploader(
            "Upload PDF di sini dan klik 'Proses'",
            accept_multiple_files=True
        )
        if st.button("Proses"):
            with st.spinner("Sedang memproses..."):
                # 1. Get PDF Text
                raw_text = get_pdf_text(pdf_docs)

                # 2. Get Text Chunks
                text_chunks = get_text_chunks(raw_text)

                # 3. Create Vector Store
                vectorstore = get_vectorstore(text_chunks)

                # 4. Create Conversation Chain
                st.session_state.conversation = get_conversation_chain(vectorstore)

                st.success("Selesai!")

if __name__ == '__main__':
    main()
\end{codeblock}

\section{Cara Menjalankan}

Buka terminal dan ketik:
\begin{codeblock}
streamlit run app.py
\end{codeblock}

Browser akan otomatis terbuka di \texttt{localhost:8501}. Upload PDF apa saja (misal manual elektronik atau makalah ilmiah), tunggu proses indexing, dan mulailah bertanya!

\chapter{Lab 2: Fine-Tuning LLM (QLoRA)}
\label{chap:lab2}

Fine-tuning adalah proses melatih ulang model agar memiliki gaya bicara atau pengetahuan spesifik. Kita akan menggunakan teknik **QLoRA** (Quantized Low-Rank Adapters) yang memungkinkan training model Llama 2 / Mistral di Google Colab gratis (T4 GPU).

\section{Persiapan Dataset}

Format data untuk fine-tuning biasanya JSONL (JSON Lines).

\begin{codeblock}[title=dataset.jsonl]
{"input": "Siapa presiden Indonesia?", "output": "Joko Widodo adalah presiden ke-7..."}
{"input": "Apa ibukota Jawa Barat?", "output": "Bandung adalah ibukota..."}
... (minimal 100-500 baris untuk hasil terlihat)
\end{codeblock}

\section{Kode Notebook (Google Colab)}

Copy kode ini ke sel Google Colab.

\subsection{1. Install Libraries}
\begin{codeblock}[language=Python]
!pip install -q -U bitsandbytes
!pip install -q -U git+https://github.com/huggingface/transformers.git
!pip install -q -U git+https://github.com/huggingface/peft.git
!pip install -q -U git+https://github.com/huggingface/accelerate.git
!pip install -q -U datasets
\end{codeblock}

\subsection{2. Load Model \& Tokenizer}
\begin{codeblock}[language=Python]
import torch
from transformers import AutoTokenizer, AutoModelForCausalLM, BitsAndBytesConfig

model_id = "bn22/Mistral-7B-Instruct-v0.1-sharded"

bnb_config = BitsAndBytesConfig(
    load_in_4bit=True,
    bnb_4bit_use_double_quant=True,
    bnb_4bit_quant_type="nf4",
    bnb_4bit_compute_dtype=torch.bfloat16
)

tokenizer = AutoTokenizer.from_pretrained(model_id)
model = AutoModelForCausalLM.from_pretrained(
    model_id,
    quantization_config=bnb_config,
    device_map={"":0}
)
\end{codeblock}

\subsection{3. Setup LoRA Config}
\begin{codeblock}[language=Python]
from peft import LoraConfig, get_peft_model

config = LoraConfig(
    r=8,
    lora_alpha=32,
    target_modules=["q_proj", "v_proj"],
    lora_dropout=0.05,
    bias="none",
    task_type="CAUSAL_LM"
)

model = get_peft_model(model, config)
\end{codeblock}

\subsection{4. Training Loop}
\begin{codeblock}[language=Python]
from transformers import TrainingArguments, Trainer

data = load_dataset("json", data_files="dataset.jsonl")
data = data.map(lambda samples: tokenizer(samples["input"]), batched=True)

trainer = Trainer(
    model=model,
    train_dataset=data["train"],
    args=TrainingArguments(
        per_device_train_batch_size=1,
        gradient_accumulation_steps=4,
        warmup_steps=2,
        max_steps=50,
        learning_rate=2e-4,
        fp16=True,
        logging_steps=1,
        output_dir="outputs",
        optim="paged_adamw_8bit"
    ),
    data_collator=transformers.DataCollatorForLanguageModeling(tokenizer, mlm=False),
)

trainer.train()
\end{codeblock}

\chapter{Lab 3: Membangun AI Agent dari Nol}
\label{chap:lab3}

AI Agent adalah sistem yang bisa menggunakan tools (Google Search, Kalkulator) untuk menyelesaikan masalah.

\section{Logika ReAct (Reason + Act)}

Kita akan menulis kode Python murni tanpa LangChain untuk memahami cara kerja Agent.

\begin{codeblock}[language=Python]
import openai
import re

class Agent:
    def __init__(self, system=""):
        self.system = system
        self.messages = []
        if self.system:
            self.messages.append({"role": "system", "content": system})

    def __call__(self, message):
        self.messages.append({"role": "user", "content": message})
        result = self.execute()
        self.messages.append({"role": "assistant", "content": result})
        return result

    def execute(self):
        completion = openai.ChatCompletion.create(
            model="gpt-4",
            messages=self.messages
        )
        return completion.choices[0].message.content

# Tools
def calculate(what):
    return eval(what)

def average_dog_weight(name):
    if name in "Scottish Terrier":
        return("Scottish Terriers average 20 lbs")
    elif name in "Border Collie":
        return("a Border Collies average weight is 37 lbs")
    elif name in "Toy Poodle":
        return("a toy poodles average weight is 7 lbs")
    else:
        return("An average dog weights 50 lbs")

known_actions = {
    "calculate": calculate,
    "average_dog_weight": average_dog_weight
}

# Prompt Engineering untuk Agent
prompt = """
Anda berjalan dalam loop Thought, Action, PAUSE, Observation.
Di akhir loop, Anda memberikan Answer.
Gunakan Thought untuk mendeskripsikan pemikiran Anda.
Gunakan Action untuk menjalankan salah satu tindakan yang tersedia: calculate, average_dog_weight.
Contoh sesi:
Question: Berapa berat gabungan anjing Terrier dan Border Collie?
Thought: Saya harus mencari berat Terrier dulu.
Action: average_dog_weight: Scottish Terrier
PAUSE
Observation: Scottish Terriers average 20 lbs
Thought: Sekarang cari berat Border Collie.
Action: average_dog_weight: Border Collie
PAUSE
Observation: a Border Collies average weight is 37 lbs
Thought: Sekarang hitung totalnya.
Action: calculate: 20 + 37
PAUSE
Observation: 57
Answer: Berat gabungannya adalah 57 lbs.
""".strip()

agent = Agent(prompt)
\end{codeblock}

\section{Menjalankan Loop}

\begin{codeblock}[language=Python]
def query(question, max_turns=5):
    i = 0
    bot = Agent(prompt)
    next_prompt = question
    while i < max_turns:
        i += 1
        result = bot(next_prompt)
        print(result)

        actions = [
            re.match(r"Action: (\w+): (.*)", line.strip())
            for line in result.split('\n')
            if re.match(r"Action: (\w+): (.*)", line.strip())
        ]

        if actions:
            action, action_input = actions[0].groups()
            if action not in known_actions:
                raise Exception("Unknown action: {}: {}".format(action, action_input))
            print(" -- running {} {}".format(action, action_input))
            observation = known_actions[action](action_input)
            print("Observation:", observation)
            next_prompt = "Observation: {}".format(observation)
        else:
            return
\end{codeblock}

Coba jalankan dengan:
\texttt{query("Berapa berat mainan Poodle dikali 2?")}

\chapter{Lab 4: Computer Vision Pipeline}
\label{chap:lab4}

Kita akan membuat skrip pendeteksi wajah real-time menggunakan webcam laptop Anda.

\section{Persiapan}
\begin{codeblock}
pip install opencv-python
\end{codeblock}

\section{Kode Pendeteksi Wajah}

\begin{codeblock}[language=Python]
import cv2

# Load pre-trained Haar Cascade classifier
face_cascade = cv2.CascadeClassifier(
    cv2.data.haarcascades + 'haarcascade_frontalface_default.xml'
)

# Buka webcam (0 biasanya ID webcam default)
cap = cv2.VideoCapture(0)

while True:
    # Baca frame demi frame
    ret, frame = cap.read()
    if not ret:
        break

    # Konversi ke grayscale (CV bekerja lebih baik di BW)
    gray = cv2.cvtColor(frame, cv2.COLOR_BGR2GRAY)

    # Deteksi wajah
    faces = face_cascade.detectMultiScale(
        gray,
        scaleFactor=1.1,
        minNeighbors=5,
        minSize=(30, 30)
    )

    # Gambar kotak di sekitar wajah
    for (x, y, w, h) in faces:
        cv2.rectangle(frame, (x, y), (x+w, y+h), (0, 255, 0), 2)
        cv2.putText(frame, 'Wajah', (x, y-10),
                   cv2.FONT_HERSHEY_SIMPLEX, 0.9, (36,255,12), 2)

    # Tampilkan hasil
    cv2.imshow('Video', frame)

    # Tekan 'q' untuk keluar
    if cv2.waitKey(1) & 0xFF == ord('q'):
        break

cap.release()
cv2.destroyAllWindows()
\end{codeblock}

\chapter{Lab 5: Voice Assistant (J.A.R.V.I.S Mini)}
\label{chap:lab5}

Menggabungkan 3 teknologi: Mendengar (Whisper), Berpikir (GPT), dan Berbicara (TTS).

\section{Arsitektur}
\begin{enumerate}
    \item Rekam audio dari mic.
    \item Transkrip audio ke teks.
    \item Kirim teks ke LLM.
    \item Konversi balasan LLM ke audio.
\end{enumerate}

\section{Kode Implementasi}

\begin{codeblock}[language=Python]
import speech_recognition as sr
from openai import OpenAI
import os
import time
from gtts import gTTS # Google Text-to-Speech (Gratis)
import playsound

client = OpenAI()

def listen():
    r = sr.Recognizer()
    with sr.Microphone() as source:
        print("Mendengarkan...")
        audio = r.listen(source)
        try:
            text = r.recognize_google(audio, language="id-ID")
            print(f"Anda: {text}")
            return text
        except:
            print("Maaf, tidak terdengar.")
            return ""

def think(text):
    response = client.chat.completions.create(
        model="gpt-3.5-turbo",
        messages=[
            {"role": "system", "content": "Jawab singkat dan padat."},
            {"role": "user", "content": text}
        ]
    )
    return response.choices[0].message.content

def speak(text):
    print(f"AI: {text}")
    tts = gTTS(text=text, lang='id')
    filename = "voice.mp3"
    tts.save(filename)
    playsound.playsound(filename)
    os.remove(filename)

def main():
    while True:
        user_input = listen()
        if "keluar" in user_input.lower():
            break
        if user_input:
            response = think(user_input)
            speak(response)

if __name__ == "__main__":
    main()
\end{codeblock}

\begin{tipbox}
Untuk suara yang lebih natural, Anda bisa mengganti \texttt{gTTS} dengan \textbf{ElevenLabs API} (berbayar tapi sangat realistis).
\end{tipbox}
